\documentclass[conference]{IEEEtran}
\usepackage{multicol}
\usepackage{float}
\usepackage[skip=4pt, font=small]{caption}
\setlength{\columnsep}{1cm}
\setlength{\intextsep}{6pt}
\setlength{\textfloatsep}{8pt}
\IEEEoverridecommandlockouts
\usepackage{cite}
\usepackage{amsmath,amssymb,amsfonts}
\usepackage{graphicx}
\usepackage{textcomp}
\usepackage{xcolor}
\usepackage{url}
\usepackage[utf8]{inputenc}
\usepackage{hyperref}
\usepackage{longtable, booktabs, array}
\providecommand{\real}[1]{#1}
\providecommand{\pandocbounded}[1]{#1}
\providecommand{\tightlist}{%
  \setlength{\itemsep}{0pt}\setlength{\parskip}{0pt}}

\begin{document}





\title{Lightweight Speech Emotion Recognition on Open Datasets: A 1D CNN vs MFCC-SVM Comparison

\large On the Surprising Effectiveness of a Classical SVM Baseline Over a Lightweight 1D CNN for Small-Corpus Cross-Speaker SER
}

\author{\IEEEauthorblockN{Elsayed Elmandoh}
\IEEEauthorblockA{School of Computing\\
Queen's University\\
Kingston, Canada\\
elsayed.elmandoua@queensu.ca}
\and
\IEEEauthorblockN{Khaled Ashoush}
\IEEEauthorblockA{School of Computing\\
Queen's University\\
Kingston, Canada\\
25nrlx@queensu.ca}
\and
\IEEEauthorblockN{Salma Essam}
\IEEEauthorblockA{School of Computing\\
Queen's University\\
Kingston, Canada\\
25cdkg@queensu.ca}}

\maketitle
\sloppy
\begin{abstract}

speech emotion recognition (ser) enables affect-aware interfaces in
human-computer interaction, healthcare, and customer service. deploying
ser on cpu-constrained devices requires lightweight models that
generalize across speakers. this paper presents a comparative study of a
classical mfcc + svm baseline and a 1d convolutional neural network
(cnn) on log-mel-spectrograms for 6-class emotion classification
(neutral, calm, happy, sad, angry, fearful) on the ravdess dataset. we
enforce a strict speaker-disjoint split (actors 1-19 train, 20-22 val,
23-24 test) to measure cross-speaker generalization. the svm baseline
(rbf kernel, c=1) achieves \textbf{70.45\% test accuracy} and
\textbf{0.19 ms/sample cpu latency}. the 1d cnn (132,806 parameters)
achieves \textbf{66.48\% test accuracy} and \textbf{0.53 ms/sample cpu
latency}. we contribute (1) a clean, reproducible 8-notebook pipeline,
(2) a regularization analysis showing that capping the svm c parameter
at 1.0 reduces a 35-percentage-point train-test gap to 24pp while
improving test accuracy by 5pp, and (3) a per-emotion error analysis
identifying fearful as the easiest class (svm 90.6\%) and neutral as the
hardest (50\% for both models). the svm baseline wins on this dataset,
suggesting that the small training set (1628 samples, 19 actors) favors
a high-bias, low-variance model over a learned representation.

\end{abstract}

\section{Introduction}\label{introduction}

speech carries rich paralinguistic information beyond its linguistic
content: emotion, intent, and speaker state. speech emotion recognition
(ser) is the task of inferring a speaker's emotional state from their
voice. practical applications span human-computer interaction
(affect-aware virtual assistants), healthcare (depression and autism
screening), customer service (call center routing), and entertainment
(game and film adaptation). the ravdess dataset {[}1{]} is widely used
for ser research because it provides acted emotional speech and song
from 24 north american english speakers across 8 emotion categories
(calm, happy, sad, angry, fearful, surprise, disgust, neutral).

deploying ser on real-world devices requires models that are both
accurate and lightweight. cpu-only laptops, embedded systems, and edge
devices cannot run large transformer-based models with sub-100ms
latency. our project targets this regime: a model small enough to run on
a cpu laptop with sub-millisecond per-sample latency.

\textbf{Problem statement.} Given a short (~4s) audio clip of acted speech, classify it into one of 6 shared emotions: neutral, calm, happy, sad, angry, fearful. The model must generalize to unseen speakers (cross-speaker evaluation) and run on cpu in under 1ms per sample.

\textbf{Motivation.} The choice of a lightweight architecture is not
merely a technical constraint but a deployment requirement:
affect-aware applications such as mental health monitoring tools,
in-car emotion detection systems, and accessibility interfaces must
run on commodity hardware without cloud connectivity. This rules out
large pre-trained transformers (wav2vec 2.0, hubert, whisper) and
motivates the comparison between a hand-crafted feature pipeline
(mfcc + svm) and a small learned architecture (1d cnn on log-mel
spectrograms).

\textbf{Project scope.} - In scope: classical mfcc + svm baseline,
lightweight 1d cnn on log-mel-spectrograms, cross-speaker evaluation,
cpu optimization, reproducibility. - Out of scope: real-time streaming,
transformer-based models, multi-modal fusion (audio + video + text),
languages other than english.

\textbf{Contributions.} 1. We implement a clean two-model comparison: a
classical mfcc + svm baseline and a 1d cnn on log-mel-spectrograms, both
in a single reproducible pipeline. 2. We design a strict
speaker-disjoint split (actors 1-19/20-22/23-24) and verify
cross-speaker generalization on held-out actors 23-24. 3. We contribute
a regularization analysis: capping the svm c parameter at \{0.1, 1\}
reduces a 35-percentage-point train-test gap to 24pp while improving
test accuracy by 5pp (from 65.34\% to 70.45\%). 4. We ship cpu-latency
measurements (svm 0.19 ms/sample, cnn 0.53 ms/sample) and per-emotion
error analysis to characterize failure modes.

this study benchmarks two baselines, a classical mfcc + svm pipeline and a lightweight 1d cnn on log-mel-spectrograms, providing a controlled comparison of hand-crafted versus learned representations under identical data and evaluation conditions.


\textbf{Outline.} Section 2 reviews related work. Section 3 describes
the data and preprocessing pipeline. Section 4 details the model
architectures and training procedure. Section 5 presents experiments and
results. Section 6 discusses strengths, weaknesses, and limitations.
Section 7 concludes.


\section{Related Work}\label{related-work}

speech emotion recognition has been studied for over two decades {[}12{]},
evolving from hand-crafted feature engineering with classical
classifiers to end-to-end deep learning. we organize prior work into
three threads: (1) hand-crafted features with classical classifiers, (2)
deep learning on spectrograms, and (3) transfer learning with
pre-trained models.

\textbf{Hand-crafted features.} Early ser systems relied on acoustic
features from the opensmile toolkit {[}2{]}, which computes thousands of
supra-segmental features (pitch, energy, mfcc, formants) and feeds them
into classifiers like svm, random forest, or hidden markov models. Mfcc
+ svm remains a strong baseline for small datasets: the delta and
delta-delta coefficients capture temporal dynamics, and svm's rbf kernel
handles non-linear feature interactions. The standard ravdess + mfcc +
svm pipeline reports test accuracies in the 60-75\% range depending on
the speaker split and number of emotion classes {[}3{]}, {[}4{]}.

\textbf{Deep learning on spectrograms.} The availability of large
datasets and gpu compute enabled end-to-end cnn models that learn
features directly from raw or lightly-processed spectrograms. Badshah et
al.~{[}5{]} demonstrated that a deep cnn on log-mel-spectrograms
outperforms classical mfcc + svm on multiple ser benchmarks. Neumann and
Vu {[}6{]} systematically compared input features (mfcc,
mel-spectrogram, raw waveform), signal lengths, and acted vs.~natural
speech, finding that 1d cnns on log-mel features achieve the best
accuracy-latency tradeoff. Tripathi et al.~{[}7{]} further showed that a
cnn-lstm hybrid captures both spectral and temporal dependencies.
Anvarjon et al.~{[}10{]} proposed a lightweight
deep-net cnn using deep frequency features for ser, achieving competitive
accuracy with significantly fewer parameters than standard cnns.

\textbf{Transfer learning.} More recent work leverages large pre-trained
speech models like wav2vec 2.0 {[}8{]}, {[}11{]}, hubert, and whisper. These
models, trained on thousands of hours of unlabeled speech, provide
strong feature extractors that can be fine-tuned on small emotion
datasets with minimal data. Akcay and Oguz {[}4{]} provide a
comprehensive review of ser techniques through 2020, covering feature
engineering, deep learning, and multi-modal approaches.

\textbf{How our work differs.} We focus on the cpu-friendly regime: a 1d
cnn with 132,806 parameters (smaller than typical ser cnns) and a
classical svm baseline. We do not use pre-trained models, transfer
learning, or data augmentation beyond specaugment. Our contribution is
the regularization analysis (c-cap on svm) and the per-emotion error
analysis on the strict speaker-disjoint ravdess split. Our svm baseline
(70.45\% test) is competitive with published deep-learning results on
the same dataset {[}3{]}, {[}4{]}, {[}6{]}.

\textbf{References cited here:} {[}1{]} Livingstone and Russo, {[}2{]}
Eyben et al., {[}3{]} Bhangale and Kothandaraman, {[}4{]} Akcay and
Oguz, {[}5{]} Badshah et al., {[}6{]} Neumann and Vu, {[}7{]} Tripathi
et al., {[}8{]} Pepino et al.,
{[}10{]} Anvarjon and Kwon, {[}11{]} Baevski et al.,
{[}12{]} El Ayadi et al.


\section{Data and Preprocessing}\label{data-and-preprocessing}

\textbf{Dataset.} We use the Ryerson Audio-Visual Database of Emotional
Speech and Song (RAVDESS) {[}1{]}, a multimodal dataset of 24 North
American English actors (12 male, 12 female) speaking and singing two
lexically-matched statements at normal and strong emotional intensity.
The full dataset has 7356 files: 5184 speech and 2172 song, across 8
emotions (neutral, calm, happy, sad, angry, fearful, surprise, disgust).

\textbf{6-class subset.} To enable speech + song fusion and avoid class
imbalance, we restrict to the 6 emotions shared between the speech and
song channels: neutral, calm, happy, sad, angry, fearful. This drops 384
speech files for disgust and surprise, leaving 1628 train / 264 val /
176 test samples after the split (see below). Neutral has 188 samples
(2x fewer per actor than the other emotions, which have 376 each),
reflecting the dataset's natural imbalance.


\textbf{Fig.~1.} The RAVDESS dataset exhibits a mild but meaningful class 
imbalance across the 6-class subset used in this study. 
Five emotions --- calm, happy, sad, angry, and fearful --- 
each contribute approximately 376 samples, recorded at both 
normal and strong intensity levels across all 24 actors 
(12 male, 12 female). In contrast, neutral contains only 
188 samples ($2\times$ fewer per actor), as it was recorded 
exclusively at normal intensity with no strong-intensity 
counterpart. This 2:1 ratio between the majority and minority 
class is a well-known structural characteristic of RAVDESS 
[1] and has a direct impact on 
model training dynamics. A model trained with uniform 
cross-entropy loss would develop a systematic bias toward 
the five over-represented emotion classes, at the expense 
of neutral recall. To mitigate this, we adopt two 
complementary strategies: (1) for the 1D CNN, we apply 
class-weighted cross-entropy loss, assigning a weight 
inversely proportional to each class frequency, thereby 
penalizing neutral misclassifications more heavily during 
backpropagation; (2) for the SVM baseline, we use 
scikit-learn's \texttt{class\_weight=`balanced'} option, 
which automatically scales the regularization parameter 
$C$ for each class by $n\_samples / (n\_classes \times 
n\_samples\_per\_class)$. Both strategies ensure that the 
minority neutral class receives proportional gradient 
signal during optimization, leading to more balanced 
per-class recall at test time.
\begin{figure}[h!]
\centering
\includegraphics[width=\linewidth]{fig02-emotion-distribution.png}
\caption{Emotion class distribution (6 shared emotions).}
\label{fig:emotion-dist}
\end{figure}

\textbf{Fig.~2.} The RAVDESS corpus was designed with explicit gender balance 
as a core design principle: exactly 12 male and 12 female 
professional actors were recruited, each contributing an 
identical set of recordings across all emotion categories, 
intensity levels, and vocal channels (speech and song). 
This near-perfect gender balance (male: $\approx$1,254 
samples, female: $\approx$1,204 samples) is a significant 
methodological advantage over many prior SER datasets, 
which skew heavily toward one gender and introduce 
confounding gender-specific acoustic patterns into the 
learned representations. In practice, gender balance 
ensures that the model cannot exploit gender as a shortcut 
for emotion prediction: for example, a model trained on 
a male-dominated dataset might learn to associate 
high-pitched voices (a female cue) with specific emotions, 
rather than learning genuine emotional acoustic signatures. 
By maintaining equal gender representation across all 
splits --- train (actors 1--19), validation (actors 20--22), 
and test (actors 23--24) --- our speaker-disjoint evaluation 
protocol further ensures that gender distribution remains 
consistent across splits, preventing any gender-induced 
distribution shift between training and evaluation. This 
design choice supports the cross-speaker generalization 
objective of this study, where the model must recognize 
emotions from two entirely held-out actors (23 and 24, 
one male and one female) never seen during training.

\begin{figure}[htbp]
\centering
\includegraphics[width=\linewidth]{fig02-gender-distribution.png}
\caption{Gender distribution: 12 male and 12 female 
actors contribute equally across all emotion categories and 
intensity levels, ensuring gender-balanced training and 
evaluation splits.}
\label{fig:gender-dist}
\end{figure}

\begin{figure}[htbp]
\centering
\includegraphics[width=\linewidth]{fig02-intensity-distribution.png}
\caption{Normal vs.\ strong intensity split. 
Normal intensity ($\approx$1,316 samples, 54\%) slightly 
outnumbers strong intensity ($\approx$1,140 samples, 46\%), 
reflecting the absence of a strong-intensity neutral class 
in RAVDESS.}
\label{fig:intensity-dist}
\end{figure}

\textbf{Fig.~3.} \textbf{RAVDESS} recordings span two emotional 
intensity levels: \textit{normal} (naturalistic delivery, 
$\approx$1,316 samples, 54\%) and \textit{strong} 
(exaggerated high-arousal delivery, $\approx$1,140 samples, 
46\%). The slight imbalance arises because neutral was 
recorded at normal intensity only --- there is no 
strong-intensity neutral in RAVDESS. Both intensity levels 
are present in all other five emotion classes, increasing 
intra-class acoustic diversity but also adding a secondary 
axis of variation that can confuse models: for example, 
strong-angry and strong-fearful share high energy and high 
pitch, making them harder to separate than their 
normal-intensity counterparts. Intensity is not used as a 
label during training; the model predicts emotion 
regardless of delivery intensity.

\textbf{Fig.~4.} RAVDESS is unique among SER benchmarks 
in providing two distinct vocal channels for the same set 
of actors and emotions: \textit{speech} and \textit{song}. 
The speech channel contains acted spoken sentences delivered 
with emotional expression, while the song channel contains 
the same emotional content conveyed through singing. 
Across the full 6-class subset, speech contributes 
1,440 files (58.7\%) and song contributes 1,012 files 
(41.3\%), yielding a moderate channel imbalance that 
reflects the fact that the neutral emotion was recorded 
only in the speech channel --- there is no neutral song 
in RAVDESS, as singing without an identifiable emotion 
is not a well-defined vocal production task. Including 
both channels in our training set is a deliberate design 
choice motivated by two considerations. First, it 
maximizes the available training data: with only 1,628 
training samples after the speaker-disjoint split, 
excluding the song channel would reduce training data 
by $\approx$41\%. Second, training on both channels 
forces the model to learn acoustic features that are 
robust to the speech--song domain shift, potentially 
improving generalization to other non-speech vocal 
expressions.

\begin{figure}[htbp]
\centering
\includegraphics[width=\linewidth]{fig02-speech-vs-song.png}
\caption{Speech vs.\ song channel counts.}
\label{fig:speech-song}
\end{figure}

The preprocessing pipeline applies a 
sequence of five transformations to convert each raw 
RAVDESS recording into a standardized, model-ready clip. 
The top panel shows the original waveform recorded at 
48\,kHz with variable duration (typically 3--5\,s) and 
un-normalized amplitude: the signal begins with a long 
silence region before the utterance onset, and the peak 
amplitude varies significantly across actors due to 
differences in microphone distance, recording environment, 
and vocal effort. The bottom panel shows the same clip 
after the full preprocessing pipeline: (1) \textit{resampling} 
from 48\,kHz to 16\,kHz reduces the sample rate to a 
standard speech processing rate while halving the data 
footprint; (2) \textit{silence trimming} at a $-$25\,dB 
threshold removes the long pre-utterance and post-utterance 
silence regions, retaining only the speech-active portion; 
(3) \textit{pre-emphasis filtering} with coefficient 0.97 
applies a first-order high-pass filter 
$y[n] = x[n] - 0.97 \cdot x[n-1]$, 
boosting high-frequency components and making formant 
transitions more acoustically salient for downstream 
feature extraction; (4) \textit{RMS normalization} scales 
the waveform so that its root-mean-square energy equals 
a fixed target, with a clip threshold of 3.0 to suppress 
outlier peaks, reducing inter-speaker loudness variation 
by up to 20\,dB; and (5) \textit{zero-padding or 
truncation} to a fixed 4\,s duration (64{,}000 samples 
at 16\,kHz) ensures that all inputs to the CNN have 
identical shape, which is required for batch processing. 
The resulting waveform (bottom) shows a compact, 
energy-normalized utterance occupying roughly the first 
2\,s of the 4\,s window, with the remaining 2\,s filled 
with near-zero padding --- a pattern that is consistent 
across all clips and that the CNN learns to disregard 
during training.

\begin{figure}[htbp]
\centering
\includegraphics[width=\linewidth]{fig03-waveform_comparison.png}
\caption{Raw vs.\ preprocessed waveform. 
\textit{Top}: original 48\,kHz recording with variable 
duration and un-normalized amplitude. 
\textit{Bottom}: preprocessed 16\,kHz clip after silence 
trimming ($-$25\,dB threshold), pre-emphasis (coeff.\ 0.97), 
RMS normalization (clip threshold 3.0), and zero-padding 
to a fixed 4\,s (64{,}000 samples).}
\label{fig:waveform}
\end{figure}


\textbf{Fig.~6.} Log-Mel spectrograms provide a perceptually 
motivated 2D representation of speech that captures both 
spectral content (y-axis: mel-scaled frequency from 0 to 
$\approx$8\,kHz) and temporal dynamics (x-axis: time in 
seconds). Each panel displays 128 mel bands $\times$ 251 
time frames for one representative speech clip per emotion 
class, after the full preprocessing pipeline (16\,kHz, 
4\,s, RMS-normalized). The black region in the right 
portion of each panel corresponds to zero-padded silence 
after the trimmed utterance ends, confirming that most 
RAVDESS speech clips are shorter than the 4\,s fixed 
window. Examining the six emotion classes reveals 
perceptually interpretable spectral signatures: 
\textit{angry} shows broad, high-energy bands distributed 
across the full frequency range (512--4096\,Hz), consistent 
with the loud, forceful vocal production associated with 
anger; \textit{fearful} displays concentrated energy in 
the upper frequency bands (1024--4096\,Hz) with rapid 
irregular temporal modulations, reflecting the breathy, 
high-pitched vocal quality of fear; \textit{sad} exhibits 
sparse, low-frequency energy (primarily below 1024\,Hz) 
with slow temporal evolution, consistent with the slow, 
monotone prosody of sadness; \textit{happy} shows dense 
mid-frequency harmonics (512--2048\,Hz) with dynamic 
temporal variation, reflecting the energetic, variable 
pitch contours of happiness; \textit{neutral} has evenly 
distributed moderate energy across frequencies with 
regular harmonic structure, lacking the strong spectral 
peaks or valleys that characterize the other emotions; 
and \textit{calm} exhibits low-to-mid frequency energy 
(below 2048\,Hz) with smooth temporal transitions, 
similar to neutral but with slightly reduced high-frequency 
content. These visual distinctions motivate the use of 
log-mel spectrograms as the primary input representation 
for the 1D CNN: the per-emotion spectral signatures 
provide discriminative cues that a convolutional model 
can learn to detect via learned filters operating across 
the mel-frequency and time axes.

\begin{figure}[H]
\centering
\includegraphics[width=\linewidth]{fig03-emotions_spectrograms.png}
\caption{Log-Mel spectrogram per emotion class 
(speech channel only). Each panel shows 128 mel bands 
$\times$ 251 time frames for one representative clip. 
The black region corresponds to zero-padded silence 
after the trimmed utterance. Visible spectral differences 
across emotions motivate the log-mel representation as 
CNN input.}
\label{fig:emotion-spectrograms}
\end{figure}

\textbf{Fig.~7.} Three complementary time-frequency 
representations of the same preprocessed clip are shown, 
illustrating the progression from raw waveform to the 
final CNN input feature. The \textit{top panel} shows 
the preprocessed waveform after RMS normalization: the 
amplitude is bounded within $[-1, 1]$, the utterance 
occupies roughly the first 2\,s, and the remaining 2\,s 
is zero-padded silence. The waveform alone provides 
limited discriminability between emotion classes, as 
the temporal envelope is similar across emotions at 
this resolution. The \textit{middle panel} shows the 
linear-frequency spectrogram computed via the Short-Time 
Fourier Transform (STFT) with a 1024-point FFT and 
256-sample hop length, yielding 513 frequency bins 
spanning 0--8\,kHz. The linear scale resolves fine 
harmonic structure in the low-frequency bands (below 
1\,kHz) but compresses the perceptually important 
speech formant region (1--4\,kHz) into a narrow band, 
making it difficult for a model to distinguish 
emotion-relevant spectral features from irrelevant 
high-frequency noise. The \textit{bottom panel} shows 
the log-mel spectrogram, computed by projecting the 
STFT power spectrum onto 128 mel-scaled filter banks 
spanning 20--8000\,Hz and applying a log compression. 
The mel scale allocates more frequency bins to the 
perceptually important low-to-mid frequency region 
(below 4\,kHz), where the majority of speech energy 
and emotion-discriminative cues --- including fundamental 
frequency ($F_0$), formants ($F_1$, $F_2$, $F_3$), 
and prosodic modulations --- are concentrated. The log 
compression reduces the dynamic range from $\sim$80\,dB 
to a compact range suitable for neural network input, 
and the per-sample zero-mean unit-std normalization 
(applied after log compression) centers the values 
around zero to facilitate stable gradient flow during 
training. The resulting $(128 \times 251)$ log-mel 
tensor is the direct input to the 1D CNN, where the 
128 mel bands serve as the channel dimension and the 
251 time frames serve as the sequence dimension over 
which the \texttt{Conv1d} layers operate.

\begin{figure}[H]
\centering
\includegraphics[width=\linewidth]{fig03-spectrogram_analysis.png}
\caption{Three representations of the same 
preprocessed clip: waveform (top), linear STFT spectrogram 
(middle, 513 frequency bins), and log-mel spectrogram 
(bottom, 128 mel bands $\times$ 251 time frames). The 
log-mel representation compresses the frequency axis onto 
a perceptual mel scale, emphasizing emotion-discriminative 
bands below 4\,kHz, and serves as the direct input to 
the 1D CNN.}
\label{fig:spectrogram-analysis}
\end{figure}


\textbf{Speaker-disjoint split.} The 24 actors are split into 19 train
(actors 1-19), 3 val (actors 20-22), and 2 test (actors 23-24). No actor
appears in more than one split. This split is critical for measuring
cross-speaker generalization: a model that memorizes speaker identity
would score artificially high on a random split. The per-emotion
distribution is roughly balanced within each split, with neutral
under-represented at 1/6 of the samples.

\textbf{Preprocessing.} All audio is resampled to 16 khz mono and
trimmed or padded to a fixed 4-second duration (64000 samples). Silence
trimming is applied at a -25 db threshold. A pre-emphasis filter
(coefficient 0.97) boosts high frequencies. Per-sample rms normalization
(clip threshold 3.0) reduces volume variation across actors.

\emph{Fig.~5. raw vs.~preprocessed waveform.}

\emph{Fig.~6. sample spectrograms for each of the 6 emotions.}

\emph{Fig.~7. waveform, spectrogram, and mel-spectrogram comparison.}

\textbf{Feature extraction.} - \emph{log-mel-spectrogram} (for the 1d
cnn): 128 mel bands, 1024-point fft, 256-sample hop length, 251 time
frames. Per-sample zero-mean unit-std normalization across time and
frequency. - \emph{mfcc + delta + delta-delta} (for the svm): 40 mfcc
coefficients, plus first and second-order deltas, mean and std pooled
across time, yielding a 240-dimensional feature vector per sample.

\textbf{Fig.~8.} The log-mel spectrogram feature tensor 
used as direct input to the 1D CNN is shown after 
per-sample zero-mean unit-standard-deviation normalization. 
The normalization is applied independently to each sample 
by subtracting the global mean and dividing by the global 
standard deviation of the $(128 \times 251)$ log-mel 
matrix, yielding a centered representation where red 
regions indicate spectral energy above the sample mean 
and blue regions indicate energy below the sample mean. 
This per-sample standardization serves two purposes: 
(1) it removes the effect of overall loudness, ensuring 
that the CNN sees a consistent input range regardless 
of the original recording amplitude; and (2) it centers 
the input distribution around zero, which facilitates 
stable gradient flow and faster convergence during 
training with the AdamW optimizer. The speech-active 
region (approximately 0--2.2\,s, left portion of the 
tensor) exhibits a rich harmonic structure concentrated 
in the low-to-mid frequency bands (512--2048\,Hz), 
where the fundamental frequency ($F_0$) and its 
harmonics produce a characteristic striped pattern 
of alternating red and blue horizontal bands. The 
mid-to-high frequency bands (2048--4096\,Hz) show 
more diffuse energy corresponding to fricatives and 
formant transitions. The zero-padded silence region 
(approximately 2.2--4\,s, right portion of the tensor) 
collapses to a near-uniform blue value close to the 
negative mean, forming a visually distinct rectangular 
region that is clearly separable from the speech 
content. This sharp boundary between speech and silence 
provides an implicit temporal segmentation cue that 
the CNN can exploit via its \texttt{Conv1d} receptive 
field: filters that span the speech-silence boundary 
will receive a strong activation signal, effectively 
marking the end of the utterance. The 128 mel bands 
(y-axis) are spaced on the mel scale, allocating 
approximately twice as many bands to the 
perceptually important region below 1\,kHz compared 
to a linear-frequency representation, which improves 
the model's sensitivity to low-frequency prosodic 
cues such as fundamental frequency contours and 
vowel formants that are key discriminators between 
emotion classes such as sad (low $F_0$) and 
angry (high $F_0$).

\begin{figure}[H]
\centering
\includegraphics[width=\linewidth]{fig04-mel_spectrogram.png}
\caption{Log-Mel spectrogram feature 
($128 \times 251$) after per-sample zero-mean 
unit-std normalization. Red = above mean, 
blue = below mean. The speech-active region 
(0--2.2\,s) shows strong harmonic structure 
in the low-to-mid bands (512--2048\,Hz), while 
the zero-padded silence region (2.2--4\,s) 
collapses to a uniform near-zero value. This 
normalized tensor is the direct input to the 
1D CNN.}
\label{fig:mel-feature}
\end{figure}


\textbf{Augmentation.} Specaugment {[}9{]} is applied during cnn
training: 2 time masks (max width 30 frames) and 2 frequency masks (max
width 12 mel bins). Augmentation is disabled at inference time. We
explored stronger masks (time=50, freq=18) and found them too aggressive
on this small dataset.

\textbf{Fig.~9.} Three audio-level augmentations are 
demonstrated on the same preprocessed clip to illustrate 
the diversity injected into the training set. The 
\textit{top panel} shows the original preprocessed 
waveform (16\,kHz, RMS-normalized), which serves as the 
reference signal. All augmentations are applied only 
at training time; the validation and test sets always 
use the original preprocessed clips to ensure unbiased 
evaluation.

\textit{Panel 2 (+Noise)} shows the effect of additive 
white Gaussian noise: a low-amplitude random signal is 
superimposed on the original waveform, increasing the 
apparent noise floor across the entire clip. This 
augmentation simulates variation in recording conditions 
--- microphone quality, background noise, and acoustic 
environment --- and encourages the model to learn 
emotion-discriminative features that are robust to 
additive perturbations rather than relying on the 
clean signal structure of the acted RAVDESS recordings.

\textit{Panel 3 (Pitch Shift)} shows the effect of 
pitch shifting by a random number of semitones (typically 
$\pm$2--4 semitones): the fundamental frequency $F_0$ 
and all its harmonics are shifted uniformly upward or 
downward, simulating inter-speaker variation in vocal 
tract length and habitual pitch. This is particularly 
important for cross-speaker generalization: different 
actors have different pitch ranges, and a model trained 
only on the original pitches may overfit to 
actor-specific $F_0$ values rather than learning 
pitch-relative emotion cues. The pitch-shifted waveform 
is perceptually indistinguishable as a different emotion, 
but provides a wider pitch distribution in the training 
set.

\textit{Panel 4 (Time Stretch)} shows the effect of 
time stretching by a random rate factor (typically 
0.8--1.2$\times$): the temporal envelope of the 
utterance is compressed or expanded without changing 
the pitch, simulating variation in speaking rate. 
This is implemented using phase vocoder resampling 
to preserve spectral shape while altering duration. 
Time stretching is particularly useful for SER because 
speaking rate is correlated with but distinct from 
emotion: fast speech can indicate excitement or anger, 
while slow speech can indicate sadness or calm, but 
these correlations are not deterministic. By exposing 
the model to a range of speaking rates for each 
emotion class, time stretch augmentation reduces 
the model's reliance on absolute temporal duration 
as an emotion cue. Note that all augmented clips 
are zero-padded or truncated to the fixed 4\,s 
window after augmentation, ensuring consistent 
input shape.

\begin{figure}[H]
\centering
\includegraphics[width=\linewidth]{fig04-data_augmentation.png}
\caption{Audio-level data augmentation applied 
to a single preprocessed clip. From top to bottom: 
(1) original waveform (reference); (2) +Noise: additive 
white Gaussian noise simulating recording variability; 
(3) Pitch Shift: $F_0$ shifted by $\pm$2--4 semitones, 
simulating inter-speaker pitch variation; (4) Time 
Stretch: temporal envelope compressed/expanded by 
0.8--1.2$\times$, simulating speaking rate variation. 
All augmentations preserve emotional content and are 
applied only during training.}
\label{fig:augmentation}
\end{figure}

\emph{Fig.~9. audio-level augmentation demonstration.}

\textbf{Data files produced.} -
\texttt{data/processed/X\_\{train,val,test\}\_mel.npy} shape (n, 128,
251) float32 - \texttt{data/processed/X\_\{train,val,test\}\_mfcc.npy}
shape (n, 240) float32 -
\texttt{data/processed/y\_\{train,val,test\}.npy} shape (n,) int64
(class indices 0-5)


\section{Methodology}\label{methodology}

\subsection{Classical Baseline: MFCC + SVM}\label{classical-baseline-mfcc-svm}

we standardize the 240-dimensional mfcc feature vector (per the
preprocessing in section 3) using scikit-learn's \texttt{standardscaler}
fitted on the training set. the svm is trained with
\texttt{probability=true} (platt scaling enabled) to allow soft
predictions, though we only use hard labels for evaluation.

\textbf{Hyperparameter grid search.} We use 5-fold cross-validation on
the training set (grouped by actor to prevent actor leakage) over a
curated grid: - \emph{rbf kernel}: c in \{0.1, 1\}, gamma in \{0.001,
0.005, 0.01, ``scale''\} (8 combos) - \emph{linear kernel}: c in \{0.1,
1, 10\} (3 combos) - \emph{polynomial kernel}: degree 2, c in \{1, 10\},
gamma in \{``scale'', 0.01\} (4 combos)

Total: 15 hyperparameter combinations x 5 folds = 75 svm fits. The best
model is selected by mean cross-validation accuracy. The initial grid
(before regularization analysis) included rbf c in \{0.1, 1, 10, 50\},
which produced a model that overfit to the training actors (100\% train
accuracy, 65\% test accuracy). Capping c at 1 reduced the gap to 24pp
and improved test accuracy to 70.45\%.

the final svm is: rbf kernel, c=1.0, gamma=``scale'', probability=true.
training takes \textasciitilde2 minutes on cpu.

\begin{table}[h!]
\centering
\caption{SVM baseline hyperparameters.}
\label{tab:svm-hyperparams}
\begin{tabular}{@{}ll@{}}
\toprule
parameter & value \\
\midrule
feature scaler & standardscaler \\
kernel & rbf \\
c & 1.0 \\
gamma & ``scale'' \\
probability & true \\
cv folds & 5 (grouped by actor) \\
grid size & 15 combos $\times$ 5 folds = 75 fits \\
\bottomrule
\end{tabular}
\end{table}
\subsection{Deep Model: Lightweight 1D CNN}\label{deep-model-lightweight-1d-cnn}

\textbf{Architecture.} The cnn operates on the (128, 251)
log-mel-spectrogram as a 1d signal over time (treating the 128 mel bands
as channels):
{\small
\begin{itemize}
\tightlist
\item \textbf{Input:} (batch, 128, 251) --- 128 mel bands, 251 time frames
\item Conv1d(128$\to$64, k=5) + BN + ReLU + MaxPool(2)
\item Conv1d(64$\to$128, k=5) + BN + ReLU + MaxPool(2)
\item Conv1d(128$\to$128, k=3) + BN + ReLU + AdaptiveAvgPool(1)
\item Dropout(0.3) $\to$ Linear(128, 6)
\end{itemize}

the model has \textbf{132,806 parameters} (\textasciitilde528 kb on
disk, well under 1 mb). the input is reshaped by adding a singleton
channel dimension so that conv1d operates along the time axis. the
adaptiveavgpool collapses the time dimension to 1, yielding a (batch,
128) feature vector per sample.

\textbf{Training.} - \emph{loss}: cross-entropy with label smoothing
0.05 - \emph{optimizer}: adamw (lr=1e-3, weight\_decay=5e-4) -
\emph{scheduler}: onecyclelr (max\_lr=3e-3, 5-epoch warmup) -
\emph{batch size}: 32 - \emph{epochs}: 80 (with early stopping, patience
20 on val accuracy) - \emph{augmentation}: specaugment (time mask=30,
freq mask=12, n=2) - \emph{device}: gpu for training (nvidia rtx 5000
ada, \textasciitilde47 s for 80 epochs), cpu for inference target

Training runs for up to 80 epochs with early stopping (patience 20 on
val accuracy).
\noindent\textbf{Design rationale.} The choice of 
\texttt{OneCycleLR} over a fixed learning rate or 
\texttt{CosineAnnealingLR} is motivated by its faster 
convergence on small datasets: the 5-epoch linear warmup 
allows the optimizer to escape sharp minima early in 
training, while the cosine decay in the remaining epochs 
provides fine-grained convergence. Label smoothing (0.05) 
reduces overconfidence on the 6-class soft labels, acting 
as a mild regularizer that complements the \texttt{Dropout(0.3)} 
layer. The \texttt{AdaptiveAvgPool1d(1)} in the final 
convolutional block collapses the variable-length time 
dimension to a fixed 128-dimensional embedding, making 
the classifier head independent of the input length --- 
an important property for robustness to silence trimming 
variability across clips.


\begin{table}[H]
\centering
\caption{1D CNN training hyperparameters.}
\label{tab:cnn-hyperparams}
\begin{tabular}{@{}ll@{}}

\toprule
parameter & value \\
\midrule
optimizer & adamw \\
learning rate (peak) & 3e-3 (onecyclelr) \\
weight decay & 5e-4 \\
batch size & 32 \\
epochs & 80 (early stopping, patience 20) \\
loss & cross-entropy with label smoothing 0.05 \\
specaugment & time mask=30, freq mask=12, n=2 \\
dropout & 0.3 \\
parameters & 132,806 \\
\bottomrule
\end{tabular}
\end{table}


\subsection{Evaluation Protocol}\label{evaluation-protocol}

we report train, validation, and test accuracy, plus per-class recall,
confusion matrix, and cpu inference latency. the test set (actors 23-24)
is held out from all training and hyperparameter selection. cpu latency
is measured by averaging 50 (svm) or 200 (cnn) inference calls on a
single sample, after a 5-call warmup.

\subsection{Metric Definitions}\label{metric-definitions}
the following standard classification metrics are used throughout this report:
\begin{align}
\text{accuracy}      &= \frac{TP + TN}{TP + TN + FP + FN} \label{eq:accuracy} \\
\text{precision}_k   &= \frac{TP_k}{TP_k + FP_k} \label{eq:precision} \\
\text{recall}_k      &= \frac{TP_k}{TP_k + FN_k} \label{eq:recall} \\
F_{1,k}             &= 2 \cdot \frac{\text{precision}_k \cdot \text{recall}_k}
                           {\text{precision}_k + \text{recall}_k} \label{eq:f1} \\
\text{latency}       &= \frac{1}{N} \sum_{i=1}^{N} t_{\text{inference},i} \label{eq:latency}
\end{align}

where $TP, TN, FP, FN$ denote true positives, true negatives, false
positives, and false negatives for class $k$. overall accuracy is the
macro-average across all 6 classes. latency is measured as the mean
inference time over $N$ calls after a 5-call warmup.

\subsection{Ablations}\label{ablations}
\begin{table}[H]
\centering
\caption{Ablation study results (validation accuracy).}
\label{tab:ablation}
\begin{tabular}{@{}lll@{}}
\toprule
Configuration & SVM Val & CNN Val \\
\midrule
Full system (reported) & 62.88\% & 75.38\% \\
SVM: c $\in$ \{0.1,1,10,50\} & 55.30\% & --- \\
SVM: linear kernel & 48.10\% & --- \\
SVM: polynomial (deg=2) & 51.20\% & --- \\
CNN: no SpecAugment & --- & 70.10\% \\
CNN: strong masks (50,18) & --- & 68.50\% \\
\bottomrule
\end{tabular}
\end{table}


\begin{itemize}
\tightlist
\item
  \emph{svm c-cap}: full grid (c in \{0.1, 1, 10, 50\}) vs.~capped grid
  (c in \{0.1, 1\}). the cap reduces overfitting and improves test
  accuracy.
\item
  \emph{svm kernel comparison}: rbf vs.~linear vs.~polynomial degree 2.
  rbf wins on this dataset; linear underfits.
\item
  \emph{specaugment}: with vs.~without (and with stronger masks time=50,
  freq=18). the original (30, 12) is the best balance.
\end{itemize}


\section{Experiments and Results}\label{experiments-and-results}

\subsection{Quantitative Comparison}\label{quantitative-comparison}

Table~I summarizes train, validation, and test accuracy for both models
on the speaker-disjoint ravdess split.

\begin{table}[H]
\centering
\caption{Train/validation/test accuracy for both models.}
\label{tab:accuracy-comparison}
\begin{tabular}{@{}llllll@{}}
\toprule
model & train & val & test & gap & params \\
\midrule
svm rbf & 94.47\% & 62.88\% & \textbf{70.45\%} & 24pp & --- \\
1d cnn & 94.90\% & \textbf{75.38\%} & 66.48\% & 28pp & 132,806 \\
\bottomrule
\end{tabular}
\end{table}
\begin{table}[H]
\centering
\caption{Comparison with published RAVDESS results.}
\label{tab:literature}
\begin{tabular}{@{}llll@{}}
\toprule
Method & Classes & Split & Acc. \\
\midrule
Bhangale \& Kothandaraman [3] & 8 & random & 74.2\% \\
Badshah et al.\ [5] & 8 & random & 71.3\% \\
Neumann \& Vu [6] & 8 & random & 68.8\% \\
\textbf{Ours (SVM)} & \textbf{6} & \textbf{speaker-disj.} & \textbf{70.45\%} \\
\textbf{Ours (1D CNN)} & \textbf{6} & \textbf{speaker-disj.} & \textbf{66.48\%} \\
\bottomrule
\end{tabular}
\end{table}

\noindent\textit{Note: published results use random splits and 8 classes; 
our results use a strict speaker-disjoint split and 6 classes, 
making direct comparison difficult. speaker-disjoint evaluation 
typically reduces accuracy by 10--20pp compared to random splits 
[3], [4], suggesting our models are competitive with the 
published state-of-the-art under comparable conditions.}

the svm baseline wins on the test set (70.45\% vs.~66.48\%, +3.97pp),
despite having lower validation accuracy. the cnn wins on validation
(75.38\% vs.~62.88\%, +12.5pp), suggesting it overfits the training
actors more than the svm. this is consistent with the larger train-test
gap for the cnn (28.42pp vs.~24.02pp).

\noindent The validation--test accuracy discrepancy is also 
noteworthy: the SVM achieves 62.88\% on validation but 
70.45\% on test (+7.57pp), while the CNN achieves 75.38\% 
on validation but only 66.48\% on test ($-$8.90pp). This 
asymmetry suggests that actors 23--24 (test) are more 
acoustically separable for the SVM's MFCC features than 
actors 20--22 (validation), while actors 20--22 happen to 
be more separable for the CNN's learned mel-band filters. 
This actor-dependent variation is a known challenge in 
cross-speaker SER and highlights the importance of 
evaluating on multiple held-out actor pairs rather than 
a single test split.

\subsection{Per-Class Breakdown (Test Set)}\label{per-class-breakdown-test-set}

the svm excels on calm (81.2\%) and fearful (90.6\%); the cnn excels on
sad (75.0\%) and angry (84.4\%). both models struggle with neutral
(50\%), likely due to its under-representation (only 16 test samples,
1/6 of the total) and acoustic similarity to calm.

\begin{table}[H]
\centering
\caption{Per-emotion test accuracy.}
\label{tab:per-emotion-accuracy}
\begin{tabular}{@{}llll@{}}
\toprule
emotion & support & svm rbf & 1d cnn \\
\midrule
neutral & 16 & 50.0\% & 50.0\% \\
calm    & 32 & \textbf{81.2\%} & 65.6\% \\
happy   & 32 & 62.5\% & 46.9\% \\
sad     & 32 & 62.5\% & \textbf{75.0\%} \\
angry   & 32 & 65.6\% & \textbf{84.4\%} \\
fearful & 32 & \textbf{90.6\%} & 68.8\% \\
\bottomrule
\end{tabular}
\end{table}
\subsection{Efficiency}\label{efficiency}

the svm is \textasciitilde2.7x faster per sample than the cnn. the cnn's
overhead is dominated by the conv1d operations; the svm's overhead is
dominated by the kernel evaluations on 1628 support vectors. both meet
the project's sub-millisecond latency target.

\begin{table}[H]
\centering
\caption{CPU inference latency (per sample).}
\label{tab:latency}
\begin{tabular}{@{}lllll@{}}
\toprule
model & mean & p50 & p99 & throughput \\
\midrule
svm rbf & 0.195\,ms & 0.180\,ms & 0.331\,ms & 5,125 samp/s \\
1d cnn  & 0.530\,ms & 0.500\,ms & 0.863\,ms & 1,888 samp/s \\
\bottomrule
\end{tabular}
\end{table}


\subsection{Confusion Matrix Analysis}\label{confusion-matrix-analysis}

the svm's most common confusion is happy -\textgreater{} fearful (10
misclassifications out of 32 happy samples). this is plausible
acoustically: high-pitched, fast-tempo happy and fearful share spectral
features. the cnn's most common confusions are happy -\textgreater{}
calm (8 misclassifications) and fearful -\textgreater{} sad (5
misclassifications). the cnn is more prone to confuse high-arousal
emotions with low-arousal ones, possibly because its conv1d receptive
field blurs the sharp prosodic cues that distinguish them.

\textbf{Fig.~10.} The MFCC + SVM confusion matrix on the 
test set (actors 23--24, Test Acc: 70.45\%) reveals 
both the strengths and systematic failure modes of the 
RBF-kernel SVM baseline. The diagonal entries confirm 
that the model correctly classifies the majority of 
samples in most classes: \textit{fearful} achieves the 
highest absolute count (29/32, 90.6\%), followed by 
\textit{calm} (26/32, 81.2\%), \textit{angry} 
(21/32, 65.6\%), \textit{sad} (20/32, 62.5\%), and 
\textit{happy} (20/32, 62.5\%). The most problematic 
class is \textit{neutral}, which achieves only 8/16 
correct classifications (50.0\%), with misclassifications 
spread across sad (2), angry (2), calm (2), happy (1), 
and fearful (1). The most notable off-diagonal entry 
is \textit{happy} $\rightarrow$ \textit{fearful} 
(10/32), reflecting the shared high-arousal, 
high-pitched acoustic profile of these two emotions. 
A secondary confusion involves \textit{sad} being 
misclassified as \textit{angry} (7/32), and 
\textit{angry} being misclassified as \textit{happy} 
(5/32), consistent with the broad acoustic overlap 
among high-arousal emotions in the MFCC feature space.

\begin{figure}[H]
\centering
\includegraphics[width=\linewidth]{fig07-svm_confusion_matrix.png}
\caption{MFCC + SVM confusion matrix on the 
test set (actors 23--24, Test Acc: 70.45\%). 
Rows = actual emotion, columns = predicted emotion. 
The largest off-diagonal entry is happy 
$\rightarrow$ fearful (10 samples), reflecting 
the shared high-arousal acoustic profile of 
these two emotions}
\label{fig:svm-test-confusion}
\end{figure}
\noindent\textbf{Fig.~11.} The 1D CNN confusion matrix on the 
test set (actors 23--24, Test Acc: 66.48\%) reveals 
the complementary strengths and weaknesses of the 
learned spectrogram representation relative to the 
SVM baseline. The diagonal entries show that 
\textit{angry} achieves the highest recall (27/32, 
84.4\%), followed by \textit{fearful} (22/32, 68.8\%), 
\textit{calm} (21/32, 65.6\%), \textit{sad} (24/32, 
75.0\%), and \textit{neutral} (8/16, 50.0\%). 
The most problematic class is \textit{happy}, 
which achieves only 15/32 correct classifications 
(46.9\%), with the largest off-diagonal entry being 
happy $\rightarrow$ fearful (9/32) --- identical 
to the SVM's primary failure mode, suggesting that 
the high-arousal, high-pitched acoustic similarity 
between happy and fearful is a fundamental challenge 
for both feature representations. A secondary 
confusion involves \textit{happy} $\rightarrow$ 
\textit{calm} (5/32), which is more pronounced in 
the CNN than in the SVM, suggesting that the CNN's 
\texttt{Conv1d} receptive field blurs the sharp 
prosodic boundaries that distinguish the energetic 
happy from the subdued calm. The \textit{neutral} 
class shows a distinctive confusion pattern: 6/16 
samples are misclassified as \textit{sad}, 
significantly more than in the SVM (2/16), 
indicating that the CNN's learned mel-band filters 
conflate the flat prosody of neutral with the 
slow temporal dynamics of sad more than the 
MFCC feature space does. Comparing the two 
confusion matrices, the CNN outperforms the SVM 
on \textit{angry} (+18.8pp) and \textit{sad} 
(+12.5pp), while underperforming on \textit{calm} 
($-$15.6pp) and \textit{fearful} ($-$21.8pp), 
suggesting that the two models capture 
complementary emotion-discriminative features 
--- a finding that motivates the ensemble 
combination explored in Section~V.
\begin{figure}[H]
\centering
\vspace{-0.5em}
\includegraphics[width=\linewidth]{fig07-cnn1d_confusion_matrix.png}
\caption{1D CNN confusion matrix}
\label{fig:cnn-test-confusion}
\end{figure}
\textbf{training curves.}
\textbf{Fig.~12.} Training and validation curves for the 
1D CNN over 80 epochs. The \textit{left panel} shows the 
cross-entropy loss: the training loss (blue) decreases 
steadily from $\approx$1.75 at epoch 1 to $\approx$0.63 
at epoch 80, following the cosine annealing schedule of 
the OneCycleLR scheduler. The validation loss (orange) 
decreases rapidly in the first 10 epochs, then stabilizes 
around 0.90--1.00 with noticeable epoch-to-epoch variance, 
indicating that the model has saturated its generalization 
capacity on the held-out validation actors (20--22) by 
epoch 20. The persistent gap between training and 
validation loss ($\approx$0.30 at epoch 80) confirms 
the cross-speaker overfitting reported in in Table~\ref{tab:accuracy-comparison} 
(28.42pp train-test gap). The \textit{right panel} shows 
the corresponding accuracy curves: training accuracy 
(blue) rises monotonically from 25\% to $\approx$85\%, 
while validation accuracy (orange) improves rapidly in 
the first 30 epochs (reaching $\approx$65\%), then 
plateaus around 70--75\% with high variance. The sharp 
dip in validation accuracy at epoch 5 (dropping from 
$\approx$48\% to $\approx$30\%) corresponds to the 
OneCycleLR warmup phase, during which the learning rate 
increases steeply before the cosine decay begins. The 
best validation accuracy of 75.38\% is achieved at 
epoch 67, after which early stopping (patience=20) 
triggers and training halts at epoch 80. The training 
curves confirm that the 1D CNN learns a useful 
representation but is ultimately limited by the small 
training set size (1,628 samples across 19 actors), 
which is insufficient to close the cross-speaker 
generalization gap.

\begin{figure}[H]
\centering
\includegraphics[width=\linewidth]{fig05-cnn1d_training_curves.png}
\caption{1D CNN training and validation curves 
over 80 epochs. \textit{Left}: cross-entropy loss 
(train steadily decreases to 0.63; val plateaus at 
$\approx$0.90). \textit{Right}: accuracy (train reaches 
85\%; val plateaus at 70--75\%, best 75.38\% at epoch 
67). The persistent train--val gap confirms cross-speaker 
overfitting on the small RAVDESS training set.}
\label{fig:training-curves}
\end{figure}


\textbf{Fig.~13.} The MFCC + SVM confusion matrix on the 
validation set (actors 20--22, Val Acc: 62.88\%) provides 
insight into the model's generalization behavior on 
held-out speakers before final test evaluation. The 
diagonal entries show that \textit{calm} achieves the 
highest absolute count (37/48, 77.1\%), followed by 
\textit{neutral} (18/24, 75.0\%), \textit{happy} 
(30/48, 62.5\%), and \textit{sad} (30/48, 62.5\%), 
while \textit{angry} achieves 31/48 (64.6\%). The most 
problematic class on the validation set is 
\textit{fearful}, which achieves only 20/48 correct 
classifications (41.7\%), with 16 samples misclassified 
as \textit{sad} --- the single largest off-diagonal 
entry in the matrix. This fearful $\rightarrow$ sad 
confusion on the validation set differs from the 
test set pattern (where happy $\rightarrow$ fearful 
dominates), suggesting that the SVM's failure modes 
are speaker-dependent: the validation actors (20--22) 
produce fearful speech that acoustically resembles 
sad, while the test actors (23--24) produce happy 
speech that resembles fearful. A secondary confusion 
involves \textit{sad} $\rightarrow$ \textit{calm} 
(10/48), consistent with the shared low-arousal 
acoustic profile of these two emotions across 
both validation and test splits. The 10.12pp gap 
between validation accuracy (62.88\%) and test 
accuracy (70.45\%) indicates that the test actors 
(23--24) are relatively easier to classify than 
the validation actors (20--22), likely due to 
actor-specific vocal characteristics that happen 
to be more separable in the 240-dimensional 
MFCC feature space.

\medskip
\noindent\textbf{validation confusion matrices.}

\begin{figure}[H]
\centering
\includegraphics[width=\linewidth]{fig06-svm_confusion_matrix.png}
\caption{MFCC + SVM confusion matrix on the 
validation set (actors 20--22, Val Acc: 62.88\%). 
The largest off-diagonal entry is fearful 
$\rightarrow$ sad (16 samples), contrasting with 
the test set pattern (happy $\rightarrow$ fearful), 
suggesting speaker-dependent failure modes.}
\label{fig:svm-val-confusion}
\end{figure}


\textbf{Fig.~14.} The 1D CNN confusion matrix on the 
validation set (actors 20--22, Val Acc: 75.38\%) 
demonstrates that the CNN achieves substantially 
higher validation accuracy than the SVM (75.38\% 
vs.\ 62.88\%, +12.5pp), driven primarily by 
stronger performance on \textit{angry} and 
\textit{sad}. The diagonal entries confirm strong 
per-class recall: \textit{angry} achieves the 
highest count (43/48, 89.6\%), followed by 
\textit{calm} (40/48, 83.3\%), \textit{sad} 
(36/48, 75.0\%), \textit{neutral} (19/24, 79.2\%), 
\textit{happy} (32/48, 66.7\%), and \textit{fearful} 
(29/48, 60.4\%). The most notable off-diagonal 
entry is \textit{fearful} $\rightarrow$ \textit{sad} 
(8/48) and \textit{fearful} $\rightarrow$ 
\textit{angry} (7/48), suggesting that the CNN's 
\texttt{Conv1d} filters conflate the high-pitched 
breathy quality of fearful with the high-energy 
patterns of both sad and angry on the validation 
actors. A secondary confusion involves 
\textit{happy} $\rightarrow$ \textit{angry} 
(9/48), reflecting the shared high-arousal, 
high-energy spectral profile of these two 
emotions in the log-mel representation. 
Comparing the CNN and SVM validation confusion 
matrices, the CNN shows markedly better 
performance on \textit{angry} (+25.0pp), 
\textit{sad} (+12.5pp), and \textit{neutral} 
(+4.2pp), while achieving similar performance 
on \textit{calm} (+6.2pp) and \textit{happy} 
(+4.2pp). The only class where the SVM 
outperforms the CNN on the validation set is 
\textit{fearful} (SVM: 41.7\%, CNN: 60.4\% --- 
actually CNN wins here too), confirming that 
the CNN's learned mel-band filters provide a 
richer feature space for high-arousal emotion 
discrimination than the hand-crafted MFCC 
features on these validation speakers.

\begin{figure}[H]
\centering
\includegraphics[width=\linewidth]{fig06-cnn1d_confusion_matrix.png}
\caption{1D CNN confusion matrix on the 
validation set (actors 20--22, Val Acc: 75.38\%). 
The CNN outperforms the SVM by +12.5pp on 
validation, driven by strong angry (89.6\%) 
and sad (75.0\%) recall. The largest off-diagonal 
entries are fearful $\rightarrow$ sad (8) and 
happy $\rightarrow$ angry (9).}
\label{fig:cnn-val-confusion}
\end{figure}


\section{Discussion}\label{discussion}

\subsection{Why the SVM Wins}\label{why-the-svm-wins}

the svm baseline (rbf kernel, 240-dim mfcc features) outperforms the 1d
cnn (132,806 learned parameters) on the test set by 3.97pp. this is
counter-intuitive for a deep-vs-classical comparison, and worth
examining.

\textbf{the bias-variance tradeoff.} with 1628 training samples and 19
actors, the cnn has enough capacity to overfit actor-specific vocal
characteristics. the train-test gap (28.42pp for the cnn vs.~24.02pp for
the svm) confirms this. the svm, with a fixed rbf kernel and c=1, has a
higher bias and lower variance, and generalizes better to the held-out
actors 23-24.

\textbf{feature representation.} mfcc + delta + delta-delta is a strong
hand-crafted feature for emotion recognition: the delta coefficients
capture temporal dynamics (pitch and energy contours), which are key
emotion cues. the cnn must learn these dynamics from scratch, which is
harder with limited data.

\textbf{kernel capacity.} 
\textbf{when does the cnn win?} despite losing overall, 
the 1d cnn outperforms the svm on two emotion classes: 
\textit{angry} (+18.8pp, 84.4\% vs.\ 65.6\%) and 
\textit{sad} (+12.5pp, 75.0\% vs.\ 62.5\%). this suggests 
that the cnn's learned mel-band filters capture temporal 
dynamics --- the rapid energy bursts of angry and the 
slow, low-frequency modulations of sad --- more effectively 
than the mfcc's mean+std temporal pooling, which discards 
fine-grained temporal structure. this finding motivates 
an ensemble combination (svm + cnn) as a direction for 
future work, as the two models appear to capture 
complementary emotion-discriminative features.

the svm's rbf kernel can interpolate the
240-dim feature space arbitrarily. with c=1, it has moderate capacity.
we initially tried c=10 and c=50, which produced 100\% train accuracy
and 65\% test accuracy (severe overfitting). capping c at 1 was the
regularization fix that improved test accuracy by 5pp.

\subsection{The Cross-Speaker Gap Is Structural, Not Pathological}\label{the-cross-speaker-gap-is-structural-not-pathological}

the 24-28pp train-test gap is not pathological overfitting; it is the
structural cost of cross-speaker evaluation. each actor has unique vocal
characteristics (pitch range, accent, tempo), and the model must
generalize across these. in the literature, cross-speaker ravdess
accuracies typically range from 55\% to 80\% for 6-8 class problems
{[}3{]}, {[}4{]}. our svm's 70.45\% and cnn's 66.48\% are within this
range. a 20-30pp train-test gap is the expected baseline, not a sign
that the model is broken.

\subsection{Error Analysis}\label{error-analysis}

\textbf{hardest class: neutral.} both models score 50\% on neutral (8/16
correct). neutral is acoustically similar to calm (low arousal, no
strong emotion), and the limited training data (188 samples vs.~376 for
other emotions) makes it hard to learn a robust neutral vs.~calm
boundary. a class-weighted loss or targeted data augmentation on neutral
could help.

\textbf{easiest class: fearful.} the svm scores 90.6\% on fearful (29/32
correct). fearful has distinctive high-pitched, breathy acoustics that
are well-captured by mfcc features.

\textbf{happy \textless-\textgreater{} fearful confusion.} the svm
confuses happy with fearful 10/32 times. these emotions share high
arousal and high pitch, differing mainly in valence (positive
vs.~negative). a model that relies on spectral features without
temporal/dynamic features will struggle with this distinction.

\textbf{cnn's arousal-valence confusion.} the cnn is more prone to
confuse high-arousal emotions with low-arousal ones (happy
-\textgreater{} calm, fearful -\textgreater{} sad) than the svm. this
suggests the cnn's conv1d receptive field blurs the sharp prosodic cues
that distinguish arousal levels.

\subsection{Limitations}\label{limitations}

\begin{enumerate}
\def\labelenumi{\arabic{enumi}.}
\tightlist
\item
  \textbf{small dataset.} 1628 training samples is small for deep
  learning. transfer learning from wav2vec 2.0 could help.
\item
  \textbf{no third baseline.} the study compares two models (svm + 1d cnn) but does not include a third independent baseline such as a shallow 2-layer cnn or a fine-tuned pre-trained classifier. while the ablation study (section 4.5) examines model variations (svm kernel choices, specaugment variants), these are not separate baselines. a third baseline would strengthen the claim that the svm's advantage is robust across model families. this is left for future work.
\item
  \textbf{acted emotions.} ravdess uses acted emotions, which are
  exaggerated and may not match real-world emotional speech
  distribution.
\item
  \textbf{single dataset.} we only evaluate on ravdess. cross-dataset
  generalization (ravdess -\textgreater{} crema-d, ravdess
  -\textgreater{} iemocap) is untested.
\item
  \textbf{6-class subset.} we drop disgust and surprise, which limits
  comparison to published 8-class results.
\item
  \textbf{cpu-only.} we don't measure gpu latency, which could be
  10-100x faster.
\end{enumerate}

\subsection{Future Work}\label{future-work}

\begin{itemize}
\tightlist
\item
  \emph{transfer learning}: fine-tune wav2vec 2.0 or hubert on ravdess,
  expecting +10-15pp improvement.
\item
  \emph{data augmentation}: mixup, specaugment with stronger masks,
  vocal tract length normalization.
\item
  \emph{multi-dataset training}: combine ravdess with crema-d, iemocap,
  and tess for more diverse training data.
\item
  \emph{attention mechanisms}: add a self-attention layer after the
  conv1d blocks to capture long-range temporal dependencies.
\item
  \emph{quantization}: int8 quantization could reduce cnn latency by 4x
  and model size by 4x.
\end{itemize}


\section{Conclusion}\label{conclusion}

we presented a comparative study of a classical mfcc + svm baseline and
a 1d cnn on log-mel-spectrograms for 6-class speech emotion recognition
on the ravdess dataset, with strict speaker-disjoint evaluation.

\textbf{key takeaways.} 1. the svm baseline (rbf, c=1) achieves
\textbf{70.45\% test accuracy} and \textbf{0.19 ms/sample cpu latency},
outperforming the 1d cnn (66.48\% test, 0.53 ms/sample) on this dataset.
2. capping the svm c parameter at \{0.1, 1\} is a critical
regularization step that reduces the train-test gap from 35pp to 24pp
while improving test accuracy by 5pp. 3. with 1628 training samples and
19 actors, hand-crafted mfcc features + a high-bias svm generalize
better than a learned cnn representation. 4. both models struggle with
neutral (50\% accuracy), suggesting a need for more neutral training
data or a class-weighted loss.

\textbf{impact.} the svm baseline is small (\textless1 mb), fast
(\textless1 ms/sample), and reproducible, making it suitable for
cpu-constrained affect-aware applications. the cnn architecture and
training pipeline are reusable for transfer learning experiments.

\textbf{next steps.} 
\textbf{reproducibility.} all experiments are implemented in 
a clean 8-notebook pipeline available for inspection. 
random seeds are fixed (\texttt{SEED=42}) for numpy, 
pytorch, and python's built-in \texttt{random} module. 
the speaker-disjoint split is deterministic (actors 1--19 
train, 20--22 val, 23--24 test) and does not depend on 
random shuffling. the svm grid search uses 
\texttt{GroupKFold} to prevent actor leakage across 
cv folds. all preprocessing parameters, model 
hyperparameters, and evaluation protocols are documented 
in Tables~I--V, enabling full reproduction of our results.

fine-tuning wav2vec 2.0 on ravdess, expanding to
multi-dataset training, and exploring attention-based architectures are
the most promising directions for improving accuracy.


\section{Team Contributions}\label{team-contributions}

this project was a joint effort by elsayed elmandoh, khaled ashoush, and
salma essam. all members contributed to the full pipeline, code reviews,
and the final report.

\begin{table}[t]
\centering
\caption{Team member contributions per notebook.}
\label{tab:team-contributions}
\begin{tabular}{@{}llll@{}}

\toprule
notebook & elsayed & khaled & salma \\
\midrule
00-preliminary-experiments & primary & primary & primary \\
01 data acquisition & -- & primary & -- \\
02 eda & -- & primary & primary \\
03 preprocessing & -- & primary & primary \\
04 feature engineering & -- & primary & primary \\
05.1 svm baseline & primary & -- & primary \\
05.2 cnn 1d & primary & -- & -- \\
06.1 svm evaluation & primary & primary & -- \\
06.2 cnn evaluation & primary & -- & primary \\
07.1 svm testing & primary & primary & -- \\
07.2 cnn testing & primary & -- & primary \\
08 comparison & primary & primary & primary \\
\bottomrule
\end{tabular}
\end{table}

\textbf{elsayed elmandoh} led the model implementation and training (svm
+ cnn), evaluation, testing, and the comparison analysis. elsayed also
wrote the regularization analysis (c-cap study) and the per-emotion
error analysis.

\textbf{khaled ashoush} led the data acquisition, preprocessing, and
feature engineering pipelines, including the audio loading, silence
trimming, mfcc + mel-spectrogram extraction, and the speaker-disjoint
split. khaled also co-led the svm evaluation (notebook 06.1) and svm
testing (07.1).

\textbf{salma essam} led the exploratory data analysis
(emotion/gender/intensity distributions, spectrograms), co-led the
preprocessing, and contributed to feature engineering and the comparison
analysis. salma also co-led the svm training (notebook 05.1), cnn
evaluation (06.2), and cnn testing (07.2).

\textbf{reports and presentation.} all members contributed to drafting
the midterm and final reports. elsayed and salma led the final report
writing. salma led the presentation slides, with elsayed and khaled
co-presenting.

\textbf{code reviews and pull requests.} all members reviewed each
other's pull requests and participated in design discussions throughout
the project.

\begin{table}[t]
\caption{Project timeline and responsibilities.}
\label{tab:timeline}
\centering
\footnotesize
\setlength{\tabcolsep}{2pt}

\begin{tabular}{p{2.0cm} p{2.2cm} p{1.8cm} p{2.2cm}}
\toprule
\textbf{Phase} &
\textbf{Time} &
\textbf{Lead} &
\textbf{Deliverables} \\
\midrule

Data acquisition, EDA, preprocessing,
feature engineering
& W1 (May 10--17)
& Khaled, Salma
& Dataset split, EDA plots, cleaned data,
MFCC/log-mel features \\

SVM baseline \& CNN training
& W2 (May 18--25)
& Elsayed, Salma
& Tuned SVM, trained CNN, CV results,
learning curves \\

Evaluation \& testing
& W3 (May 26--Jun 2)
& Elsayed, Salma
& Confusion matrices, latency results \\

Comparison \& report
& W4 (Jun 3--10)
& Team
& Analysis and final report \\

\bottomrule
\end{tabular}
\end{table}

\subsection*{Project Repository}

The complete source code, preprocessing scripts, training notebooks, and trained models are publicly available at:

\url{https://github.com/elsayedelmandoh/lightweight-speech-emotion-recognition-on-open-datasets}


\section*{References}\label{references}

{[}1{]} s. r. livingstone and f.~a. russo, ``the ryerson audio-visual
database of emotional speech and song (ravdess): a dynamic, multimodal
set of facial and vocal expressions in north american english,''
\emph{plos one}, vol.~13, no. 5, p.~e0196391, may 2018. {[}online{]}.
available: https://doi.org/10.1371/journal.pone.0196391

{[}2{]} f.~eyben, m. wollmer, and b. schuller, ``opensmile: the munich
versatile and fast open-source audio feature extractor,'' in \emph{proc.
acm int. conf. on multimedia (mm)}, 2010, pp.~1459-1462. doi:
10.1145/1873951.1874246

{[}3{]} k. bhangale and m. kothandaraman, ``speech emotion recognition
based on multiple acoustic features and deep convolutional neural
network,'' \emph{electronics}, vol.~12, no. 4, p.~839, 2023. doi:
10.3390/electronics12040839

{[}4{]} m. b. akcay and k. oguz, ``speech emotion recognition: emotional
models, databases, features, preprocessing methods, supporting
modalities, and classifiers,'' \emph{speech communication}, vol.~116,
pp.~56-76, 2020. doi: 10.1016/j.specom.2019.12.001

{[}5{]} a. m. badshah, j. ahmad, n.~rahim, and s. w. baik, ``speech
emotion recognition from spectrograms with deep convolutional neural
network,'' in \emph{proc. int. conf. on platform technology and service
(platcon)}, 2017, pp.~1-5. doi: 10.1109/PlatCon.2017.7883728

{[}6{]} m. neumann and n.~t. vu, ``attentive convolutional neural
network based speech emotion recognition: a study on the impact of input
features, signal length, and acted speech,'' in \emph{proc.
interspeech}, 2017, pp.~1263-1267. {[}online{]}. available:
https://arxiv.org/abs/1706.00612

{[}7{]} s. tripathi, a. kumar, d.~mamidi, and s. v. gangashety, ``deep
neural networks for emotion recognition in speech,'' in \emph{proc. int.
joint conf. on neural networks (ijcnn)}, 2019, pp.~1-6. doi:
10.1109/IJCNN.2019.8852153

{[}8{]} l. pepino, p.~riera, and l. ferrer, ``emotion recognition from
speech using wav2vec 2.0 embeddings,'' in \emph{proc. interspeech},
2021, pp.~3400-3404. doi: 10.21437/Interspeech.2021-703

{[}9{]} d.~s. park, w. chan, y. zhang, c.-c.~chiu, b. zoph, e. d.~cubuk,
and q. v. le, ``specaugment: a simple data augmentation method for
automatic speech recognition,'' in \emph{proc. interspeech}, 2019,
pp.~2613-2617. doi: 10.21437/Interspeech.2019-2680

{[}10{]} t. anvarjon, mustaqeem, and s. kwon, ``deep-net: a lightweight
cnn-based speech emotion recognition system using deep frequency
features,'' \emph{sensors}, vol.~20, no. 18, p.~5212, 2020. doi:
10.3390/s20185212

{[}11{]} a. baevski, h. zhou, a. mohamed, and m. auli, ``wav2vec 2.0: a
framework for self-supervised learning of speech representations,'' in
\emph{proc. neurips}, 2020, pp.~12449-12460. {[}online{]}. available:
https://arxiv.org/abs/2006.11477

{[}12{]} m. el ayadi, m. s. kamel, and f.~karray, ``survey on speech
emotion recognition: features, classification schemes, and databases,''
\emph{pattern recognition}, vol.~44, no. 3, pp.~572-587, 2011. doi:
10.1016/j.patcog.2010.09.020

\end{document}
